% Generated draft: requires mathematical and source review.
\documentclass[11pt]{article}
\usepackage[margin=25mm]{geometry}
\usepackage{fontspec}
\setmainfont{Latin Modern Roman}
\usepackage{amsmath,amssymb,mathtools,unicode-math}
\setmathfont{Latin Modern Math}
\usepackage{enumitem,needspace,xurl,hyperref}
\setlength{\parindent}{0pt}
\setlength{\parskip}{5pt}
\setlength{\emergencystretch}{3em}
\newcommand{\N}{\mathbb{N}}
\newcommand{\Z}{\mathbb{Z}}
\newcommand{\Q}{\mathbb{Q}}
\newcommand{\R}{\mathbb{R}}
\newcommand{\C}{\mathbb{C}}
\newcommand{\F}{\mathbb{F}}
\newcommand{\A}{\mathbb{A}}
\newcommand{\PP}{\mathbb{P}}
\newcommand{\E}{\mathbb{E}}
\newcommand{\eps}{\varepsilon}
\newcommand{\dd}{\,\mathrm{d}}
\begin{document}
\section*{UnsolvedMath ID: 1100112}
\textbf{Source problem code:} catalogue entry
\subsection{Definitions and mathematical statement}
A finitely generated group $\Gamma$ is hyperbolic if, for a finite generating set, its Cayley graph with the word metric has $\delta$-thin geodesic triangles for some $\delta\geq0$. It is one-ended if the complement of every finite vertex set in this graph has exactly one infinite connected component. For $L\leq\Gamma$, the index $[\Gamma:L]$ is the number of left cosets of $L$.

\textbf{Question (Whyte).} Does there exist a one-ended hyperbolic group $\Gamma$ and subgroups $H,K\leq\Gamma$ such that
\[
[\Gamma:H]<\infty,\qquad [\Gamma:K]=\infty,\qquad H\cong K?
\]
Thus the question concerns the possibility of such a group, not an assertion that every one-ended hyperbolic group has such subgroups. The supplied research summary reports an affirmative answer in the literature and attributes an example to Stark--Woodhouse; this is the record's reported status, not a proof supplied here.
\subsection{Short English statement}
Can one one-ended hyperbolic group contain isomorphic subgroups of finite and infinite index? The record reports that such examples exist.
\subsection{Sources}
\begin{itemize}
\item M. Bestvina, \emph{Questions in Geometric Group Theory} (2004), Question 1.12, PDF page 3, attributed there to Whyte: \url{https://www.math.utah.edu/~bestvina/eprints/questions-updated.pdf}. The record describes manual transcription and access on 2026-07-29.
\item The supplied research summary attributes an affirmative example to Stark--Woodhouse, section 2, without a title, year, or bibliographic identifier. Its sketch also cites Bestvina--Feighn for hyperbolicity and Neumann's covering lemma; no further references are supplied.
\item Local registry \texttt{data/unsolvedmath/problems.json}, permanent ID $1100112$, source code catalogue entry, \emph{Questions in Geometric Group Theory --- Q 1.12}. The record labels this solved and SOLVED-IN-LITERATURE; the source-list status and rights checks remain marked NEEDS\_REVIEW.
\end{itemize}
\end{document}
